\documentclass[11pt]{article}

\usepackage[margin=1in]{geometry}
\usepackage{booktabs}
\usepackage{longtable}
\usepackage{array}
\usepackage{enumitem}
\usepackage{ragged2e}
\usepackage{xcolor}
\usepackage[hidelinks]{hyperref}
\usepackage{titlesec}
\usepackage{parskip}
\usepackage{framed}

\newcolumntype{L}[1]{>{\RaggedRight\arraybackslash}p{#1}}

\titleformat{\section}{\large\bfseries}{\thesection.}{0.5em}{}
\titleformat{\subsection}{\normalsize\bfseries}{}{0em}{}

%% ---------------------------------------------------------------------------
%% Drafting aids. REMOVE \draftmode BEFORE SUBMISSION: it turns off the banner
%% and makes any remaining \todo render loudly so nothing slips through.
%% ---------------------------------------------------------------------------
\newif\ifdraftmode\draftmodetrue

\newcommand{\todo}[1]{\textcolor{red}{\textbf{[TODO: #1]}}}
\newcommand{\status}[1]{\textcolor{gray}{\small\textsf{[#1]}}}

%% Reviewer text, quoted verbatim.
\newenvironment{rev}{\begin{leftbar}\itshape\small}{\end{leftbar}}

\title{\vspace{-1.5em}Response to Reviewers}
\author{Spatial Modeling of Tissues for Morphogenic Field Analysis\\
\small PLOS ONE, PONE-D-26-04349 \quad (Osgood-Zimmerman \& Raredon)}
\date{Draft of 2026-07-27 \quad$\vert$\quad Revision deadline: 2026-08-07}

\begin{document}
\maketitle
\vspace{-1em}

\ifdraftmode
\begin{center}
\fcolorbox{red}{red!5}{\parbox{0.95\textwidth}{\centering
\textcolor{red}{\textbf{INTERNAL DRAFT, NOT FOR SUBMISSION.}}\\[2pt]
\small Items marked \textcolor{red}{\textbf{[TODO: \ldots]}} are not yet done. Statuses: \status{DONE} the work is
finished and verified; \status{DRAFTED} response text is written but the corresponding
manuscript edit has not been made; \status{TODO} not started. Set
\texttt{\textbackslash draftmodefalse} and confirm no \texttt{\textbackslash todo}
remains before submitting.}}
\end{center}
\vspace{0.5em}
\fi

\noindent Dear Dr.\ D'Acquisto and Reviewers,

\medskip

\noindent Thank you for the careful and constructive reviews. We are grateful for the
time invested, and we think the manuscript is substantially better for the revision.

Reviewer 2 asked whether the framework identifies biologically meaningful spatial
structure beyond smoothing, increased feature dimensionality, or clustering. We regarded
this as the central question of the review, and we have addressed it directly with three
new controls (Section R2.8): a ligand-receptor pairing shuffle, a coordinate shuffle with
full model refitting, and a concordance test against the independent anatomical
annotations published with the source dataset. The coordinate-shuffle control is the most
informative of the three, and it is unambiguous: refitting the identical pipeline on
spatially permuted data collapses the emergent domains entirely.

We have also moderated the generalizability and biological-interpretation language
throughout, clarified the statistical model and its implementation, and added a
supplementary table mapping every mathematical object in the manuscript to the code that
implements it.

Reviewer comments are reproduced verbatim in italics, followed by our response. Changes
are marked in the accompanying \emph{Revised Manuscript with Track Changes}.

\medskip
\noindent Sincerely,\\
Aaron Osgood-Zimmerman and Micha Sam Brickman Raredon

\vspace{1em}
\hrule

\section{Status summary}

\ifdraftmode
\noindent\textcolor{red}{\textbf{Internal only, delete this section before submission.}}
\medskip
\fi

\begin{longtable}{@{}L{1.6cm} L{9.2cm} L{3.6cm}@{}}
\toprule
\textbf{Item} & \textbf{Ask} & \textbf{Status} \\
\midrule
\endhead
JR-1  & PLOS style and file naming & Done \\
JR-2  & Role of funder statement & Needs author input \\
JR-3  & Remove funding text from Acknowledgments & To do \\
JR-4  & PLOS \LaTeX{} template & Done \\
JR-5  & Upload Supplemental Materials as SI & To do \\
JR-6  & Evaluate recommended citations & Drafted (see R2.10) \\
\midrule
R1.1  & Comparative analysis vs existing methods & Drafted (see R2.8, R2.10) \\
R1.2  & Tissue sample and methods detail & To do \\
R1.3  & Figures too small, reorganize & To do \\
R1.4  & Small-scale disease proof of concept & Drafted (respectful pushback) \\
R1.5  & Better graphical abstract & To do \\
\midrule
R2.1  & Define model parameters and priors at first use & Drafted \\
R2.2  & Notation-to-code supplementary table & Drafted, table not built \\
R2.3  & Clarify that final analyses used observed nUMI & Drafted \\
R2.4  & Define metrics; Table 1 abbreviations and percentages & To do \\
R2.5  & Rationale for Cd44 as representative feature & Needs author input \\
R2.6  & Elongated-pixel distortion in Figs 3 and 4 & To do \\
R2.7  & Clearer captions for Figs 5 and 6 & To do \\
\textbf{R2.8} & \textbf{Controls and baselines} & \textbf{Done, all three} \\
R2.9  & Reproducibility documentation & Partly done \\
R2.10 & Positioning vs related methods & Drafted \\
R2.11 & Moderate generalizability claims & Drafted \\
R2.12 & Cautious biological language & Drafted \\
\bottomrule
\end{longtable}

\section{Journal requirements}

\subsection{JR-1. PLOS ONE style requirements and file naming} \status{DONE}

\begin{rev}
Please ensure that your manuscript meets PLOS ONE's style requirements, including those
for file naming.
\end{rev}

\noindent The manuscript has been reformatted to the PLOS ONE style templates. Section
headings, title capitalization, the reference style, and file naming now follow the PLOS
ONE guidelines. See also JR-4.

\subsection{JR-2. Role of the funder} \status{NEEDS AUTHOR INPUT}

\begin{rev}
Please state what role the funders took in the study. If the funders had no role, please
state: ``The funders had no role in study design, data collection and analysis, decision
to publish, or preparation of the manuscript.''
\end{rev}

\noindent We confirm the following statement, which we have included in our cover letter
for the journal to enter into the online submission form:

\begin{quote}
The funders had no role in study design, data collection and analysis, decision to
publish, or preparation of the manuscript.
\end{quote}

\todo{Both authors must confirm this is accurate for T32GM086287 (NIGMS) and for the Yale
and Bucknell startup funds before it goes in the cover letter.}

\subsection{JR-3. Funding text in the Acknowledgments} \status{TODO}

\begin{rev}
We note that you have provided funding information that is not currently declared in your
Funding Statement. However, funding information should not appear in the Acknowledgments
section or other areas of your manuscript.
\end{rev}

\noindent The funding text has been removed from the Acknowledgments section of the
manuscript. The Acknowledgments now contain only non-financial acknowledgments. The
funding information appears solely in the Funding Statement, and our updated statement is
reproduced in the cover letter.

\todo{Actually remove the funding sentence from the Acknowledgments in paper-plos.tex,
around line 376, and check that no funding language survives elsewhere in the text.}

\subsection{JR-4. PLOS \LaTeX{} template} \status{DONE}

\begin{rev}
Please update your submission to use the PLOS LaTeX template.
\end{rev}

\noindent The manuscript has been ported from the Elsevier \texttt{elsarticle} class to
the PLOS \LaTeX{} template, including the PLOS preamble, title and author block, and the
\texttt{plos2025} bibliography style.

\subsection{JR-5. Supporting Information upload} \status{TODO}

\begin{rev}
Please upload a copy of Supporting Information Figure/Table/etc.\ ``Supplemental
Materials'' which you refer to in your text on page 19.
\end{rev}

\noindent The supplemental materials, previously available only through external
repository links, are now uploaded as Supporting Information files S1 through S8 and are
cited individually at the appropriate points in the text.

\todo{Package the eight Supplements for Publication directories as PLOS SI files, confirm
each S-number in the manuscript matches the uploaded file, and check the SI captions
render.}

\subsection{JR-6. Recommended citations} \status{DRAFTED}

\begin{rev}
If the reviewer comments include a recommendation to cite specific previously published
works, please review and evaluate these publications to determine whether they are
relevant and should be cited.
\end{rev}

\noindent We evaluated every method named by Reviewer 2 (SpatialDE, BOOST-GP, Spatial
Variance Component Analysis, and SpatialDM) and judged all of them relevant. Each is now
cited and discussed in a new positioning subsection; see our response to R2.10. We have
additionally cited CytoSignal (Liu et al., \emph{Nature Genetics}, 2026), concurrent work
that was not available at the time of our original submission and that analyzes the same
source dataset.

\section{Academic editor}

\begin{rev}
It seems to me that the Reviewers have asked a number of clarifications and more detailed
explanation for the approaches used in the study while finding the overall idea of
interest. I would suggest the authors to fully address the comments made by the Reviewers
as this might significantly improve the impact of the study and its overall take-home
message.
\end{rev}

\noindent We agree, and we have treated every reviewer comment as actionable. The
clarifications requested by Reviewer 2 concerning the statistical model, its
implementation, and the evaluation metrics were fair; the manuscript understated how much
detail a reader needs to reproduce or critically assess the work. The point-by-point
responses below describe each change.

\section{Reviewer 1}

\subsection{R1.1. Comparative analysis with existing spatial modeling techniques} \status{DRAFTED}

\begin{rev}
Recommend the authors to perform a comparative analysis with other existing techniques
for spatial modelling techniques as verifications.
\end{rev}

\noindent We have added comparison and validation along three axes, described in full in
our response to Reviewer 2's comment R2.8.

First, within our own framework we compare candidate convolution operators
quantitatively, and this comparison now includes the product operator used by CytoSignal.
The common-minimum operator we adopt yields more spatially coherent domains than the
product operator on this dataset.

Second, we compare the emergent domains against an external reference that we did not
produce: the anatomical annotations published with the source dataset by Chen et al.
(2022). This is a comparison to an independently derived, expression-based segmentation
of the same tissue section.

Third, we added shuffle-based null controls, which are the standard validation device in
this literature.

We want to be straightforward about what we have not done. We have not run a full
method-versus-method benchmark against SpatialDE, BOOST-GP, SVCA, SpatialDM, or
CytoSignal. Those methods answer different questions from ours: they test individual
features or interactions for spatial significance, whereas our contribution is a
continuous-field representation with quantified uncertainty and the clustering of the
full high-dimensional signaling field. There is no shared target quantity on which a
head-to-head accuracy comparison would be meaningful without first defining a benchmark
task, which we regard as its own piece of work. We have instead positioned our approach
against each of these methods conceptually and explicitly (R2.10), and we now state this
limitation in the Discussion rather than leaving it implicit.

\subsection{R1.2. Detail on the tissue sample and methods} \status{TODO}

\begin{rev}
The methods on the tissue/patient sample obtained here is missing and lacks details.
\end{rev}

\noindent We have expanded the dataset description in Materials and methods. The revised
text specifies the source publication and accession, the developmental stage and section
identity, the platform and its capture chemistry, the physical dimensions of the field of
view analyzed, the binning applied, the number of spatial locations and features modeled,
and the filtering that produced the analyzed feature set.

\todo{Write this expansion into the Dataset selection and Dataset preprocessing
subsections. Concrete facts to state: Chen et al.\ 2022 Stereo-seq mouse embryo (MOSTA),
section E16.5\_E1S3; field of view x from 15000 to 25000 and y from 11500 to 18000 in raw
chip coordinates; bin size 50; 26000 spatial locations; 918 curated features of which 912
produced converged fits; FANTOM5 ligand-receptor pairings, of which 1477 have both
partners among the converged features. Note that the sample is a mouse embryo from a
public atlas, not a patient sample, which is what the reviewer may be assuming.}

\subsection{R1.3. Figure size and organization} \status{TODO}

\begin{rev}
The data presented in the figure is on the small side of the spectrum, recommend the
authors to reorganise them.
\end{rev}

\noindent We have regenerated the figures at larger physical size and higher resolution,
and reorganized several panels so that the comparisons a reader is asked to make are
adjacent. Figures 3 and 4 have also been corrected for the aspect-ratio distortion noted
by Reviewer 2 (R2.6).

\todo{Regenerate all figures. Coordinate with R2.6 (aspect ratio) and R2.7 (captions) so
each figure is touched once.}

\subsection{R1.4. Small-scale disease proof of concept} \status{DRAFTED, PUSHBACK}

\begin{rev}
Since the author claims the spatial modelling provides multidimensional information to
guide disease assessment, it might be good for the author to use this paper to perform a
small scale POC work.
\end{rev}

\noindent We appreciate this suggestion and have thought about it carefully, but we have
concluded that adding a disease dataset is not the right change for this manuscript, and
we have instead removed the framing that invited it.

The manuscript did not report a disease result; disease was raised only as future
motivation. On reflection, that framing promised more than the work delivers, which is
also the substance of Reviewer 2's comment R2.11 about generalizability. Rather than
adding a second dataset, we have moderated the disease language so the stated scope
matches the evidence: this is a methodological proof of concept on healthy developing
tissue.

We think this is the more honest resolution. A disease proof of concept done properly
requires matched healthy and diseased tissue, attention to the confounding between disease
state and tissue composition, and enough replication to distinguish a disease effect from
between-sample variation. Done in the time available on a single additional section, it
would produce an illustrative figure that we could not defend statistically, which is the
opposite of what this manuscript argues for. We have noted the application to diseased
tissue as clearly future work in the Discussion, alongside the lung-tissue direction
already mentioned there.

\subsection{R1.5. Graphical abstract} \status{TODO}

\begin{rev}
Minor comment: The graphical abstract submitted did not do the work justice, consider a
better images rather than copy-pasting from the data section.
\end{rev}

\noindent We agree, and we have replaced the graphical abstract with a purpose-designed
schematic rather than panels reused from the figures.

\todo{Design and produce the new graphical abstract. Raredon.}

\section{Reviewer 2}

\noindent We thank Reviewer 2 for an unusually thorough and useful review. The comments
identified real gaps, particularly the absence of controls, and the manuscript is
materially stronger as a result.

\subsection{R2.1. Precision of the statistical model specification} \status{DRAFTED}

\begin{rev}
First, the statistical model is not specified with sufficient precision. The manuscript
introduces a hierarchical statistical framework, but several key parameters are not
clearly defined, including $\theta_i$, $\theta_N$, $\Sigma(\theta_i)$, $\Sigma(\theta_N)$,
the covariance function, prior structure, and the relationship between the abstract
mathematical model and the computational implementation. Since this model is central to
the manuscript, all parameters should be introduced clearly at first use.
\end{rev}

\noindent We accept this criticism. The model section was written at a level of
abstraction that made the framework readable but left the actual fitted model
underdetermined. We have rewritten it to define every symbol at first use and to state
the concrete choices used in the reported analyses.

Specifically, the revised Materials and methods now state that the spatial random effect
for feature $i$ is a zero-mean Gaussian process with a Mat\'ern covariance function,
represented through the stochastic partial differential equation approach, so that
$\theta_i = (\rho_i, \sigma_i)$ collects the range and marginal standard deviation for
feature $i$, and $\Sigma(\theta_i)$ denotes the resulting Mat\'ern covariance matrix
evaluated at the mesh nodes. The smoothness parameter is fixed rather than estimated
($\alpha = 2$, giving Mat\'ern smoothness $\nu = 1$ in two dimensions). We place
penalized-complexity priors on $\theta_i$: the range prior satisfies
$P(\rho_i < 500) = 0.95$ and the marginal standard deviation prior satisfies
$P(\sigma_i > 0.1) = 0.05$, with the range in the same physical units as the chip
coordinates. A sum-to-zero constraint is applied to the field for identifiability against
the intercept. The symbols $\theta_N$ and $\Sigma(\theta_N)$ played a role only in the
latent total-count model, which the final analyses did not use; see R2.3, where we have
clarified their status rather than leaving them ambiguous.

We have also added an explicit paragraph connecting the abstract model to its
implementation, and a supplementary table making the mapping object by object (R2.2).

\todo{Write this into the ``A general spatial statistical model for spatialomics''
subsection of paper-plos.tex. The numerical prior values above are read directly from the
fitting code (\texttt{inla.spde2.pcmatern} with \texttt{prior.range = c(500, 0.95)},
\texttt{prior.sigma = c(0.1, 0.05)}, \texttt{alpha = 2}, \texttt{constr = TRUE}); confirm
they match the canonical 2025-04-10 run before the statement goes in the paper.}

\subsection{R2.2. Notation-to-code mapping} \status{DRAFTED, TABLE NOT BUILT}

\begin{rev}
Second, the correspondence between the notation in the manuscript and the public code
repository is difficult to follow. Readers should be able to identify where mathematical
objects such as $\lambda_i(s)$, $F_i(s)$, $N(s)$, $\alpha_i$, $\Sigma(\theta_i)$, and
$\theta_i$ are implemented in the code. I recommend adding a supplementary table mapping
each manuscript term to its corresponding code object, script, or model component.
\end{rev}

\noindent We have added this table as Supporting Information. Each row gives a
mathematical symbol, its meaning, the script in which it is constructed, and the name of
the corresponding code object, so that a reader can move from any equation in the
manuscript to the lines that implement it.

\todo{Build the table. One row per symbol, minimally: $\lambda_i(s)$, $F_i(s)$, $N(s)$,
$\alpha_i$, $\theta_i$, $\Sigma(\theta_i)$, $\theta_N$, $\Sigma(\theta_N)$, the mesh, the
Mat\'ern SPDE object, the likelihood, and the posterior predictive draws. Anchor each row
to a script and object name in \texttt{Supplements for Publication/code/spatial-modeling/}.
Verify every path and object name against the actual files rather than from memory.}

\subsection{R2.3. Treatment of total counts $N(s)$} \status{DRAFTED}

\begin{rev}
Third, the treatment of total counts $N(s)$ requires clarification. The manuscript
presents a latent model for $N(s)$, but later states that modeling $N(s)$ did not
substantially alter the results and that observed nUMI values were used directly. The
authors should clearly state whether the final reported analyses used a joint latent
model, an observed offset or exposure term, or a simplified model using observed nUMI
values.
\end{rev}

\noindent The reviewer is right that this was ambiguous, and we have stated it plainly.

All analyses reported in the figures and in Table 1 use the observed number of unique
molecular identifiers at each location as a fixed offset, that is, as an exposure term
entering the Poisson mean multiplicatively and treated as known rather than modeled. No
result in the manuscript comes from the joint latent model for $N(s)$.

We retain the latent formulation in the model development section because it is the
general form of the framework and because we tested it, but we now flag its status
explicitly at first presentation, rather than only in passing later. We also state what
the comparison showed: fitting the joint latent model for $N(s)$ did not materially change
the estimated feature fields, which is why the simpler and computationally cheaper offset
specification was adopted for the full feature set. The supporting comparison is the
analysis in Supplement 3.

\todo{Add the clarifying sentence at the point where the latent $N(s)$ model is first
introduced, not only in the Results. Confirm the Supplement 3 comparison supports the
``did not materially change'' wording, and quantify it if the numbers are available.}

\subsection{R2.4. Prediction metrics and model comparison} \status{TODO}

\begin{rev}
Fourth, the prediction metrics and model comparison need clearer explanation. The
manuscript compares Poisson, ZIP, and ZAP models using MAE, RMSE, Dawid-Sebastiani score,
and log score, but the definitions, orientation, and interpretation of these metrics are
not sufficiently clear. Table 1 should define all abbreviations and explain the
percentages in parentheses. The authors should also clarify which criteria were
prioritized and why the evidence supports selecting the Poisson model over ZIP and ZAP.
\end{rev}

\noindent We have added explicit definitions for each metric, stated the orientation of
each, expanded the Table 1 caption to define every abbreviation and the parenthetical
percentages, and stated which criteria we prioritized and why.

In summary: mean absolute error (MAE) and root mean squared error (RMSE) are point-forecast
accuracy measures; the Dawid-Sebastiani score (DS) and logarithmic score (LS) are proper
scoring rules that assess the full predictive distribution rather than a point summary. All
four are negatively oriented, so smaller values indicate better performance. We prioritized
the proper scoring rules, because the manuscript's argument concerns the quality of the
predictive distribution and its uncertainty rather than point accuracy alone.

\todo{This response is a skeleton and needs the actual justification for Poisson over ZIP
and ZAP written against the Table 1 numbers, including the case where the metrics disagree.
Also state exactly what the parenthetical percentages in Table 1 are (percentage of
features for which that model won on that metric, if that is what they are; verify against
the generating code). Write the metric definitions as equations in the Methods.}

\subsection{R2.5. Rationale for Cd44 as the representative feature} \status{NEEDS AUTHOR INPUT}

\begin{rev}
For example, the rationale for selecting Cd44 as a representative example should be
explained.
\end{rev}

\noindent We have added the rationale to the text and to the figure caption.

\todo{Needs the authors' actual reason. Candidate justifications, to be confirmed rather
than assumed: Cd44 sits at a typical expression level and sparsity for the modeled feature
set, so it is representative rather than favourable; it is broadly expressed across several
anatomical regions, so the raw-versus-normalized contrast is visible; it is a well
characterized surface receptor. State the real reason, and if the feature was chosen for
visual clarity, say so and note that the full per-feature results are in the supplement.}

\subsection{R2.6. Elongated-pixel distortion in Figures 3 and 4} \status{TODO}

\begin{rev}
Figures 3 and 4 appear visually distorted due to elongated rectangular pixels, which
should be corrected or clearly identified as a visualization artifact.
\end{rev}

\noindent We have regenerated Figures 3 and 4 with a fixed one-to-one aspect ratio, so
that spatial bins render as squares and distances are visually comparable along both axes.

\todo{Regenerate with a fixed coordinate ratio. Check whether the distortion originated in
the plotting call or in the underlying bin geometry; if the bins genuinely are not square,
say so explicitly instead of rescaling the display.}

\subsection{R2.7. Captions for Figures 5 and 6} \status{TODO}

\begin{rev}
Figures 5 and 6 require clearer captions, definitions of abbreviations, explanation of
rows and columns, and a more explicit connection between the visual evidence and the
conclusions drawn.
\end{rev}

\noindent We have rewritten both captions to be self-contained: each now defines every
abbreviation used in the panel, states what the rows and columns represent, and says
explicitly what the reader should take from the figure and how it supports the claim made
in the text.

\todo{Rewrite both captions. Figure 5 is the convolution-operator comparison (product,
geometric mean, common minimum, with ELSA entrograms) and Figure 6 is the feature-depth
sweep producing emergent domains. Each caption must be readable without the surrounding
prose.}

\subsection{R2.8. Controls and baselines} \status{DONE}

\begin{rev}
The manuscript would also benefit from additional controls or comparative analyses. At
present, there is no clear baseline against which the proposed framework can be evaluated.
Possible additions include comparison with existing spatial transcriptomics or
ligand-receptor methods, simulated datasets with known ground truth, spatial shuffling
controls, ligand-receptor pair shuffling controls, random ligand-receptor combinations, or
comparison with previously annotated tissue regions. These analyses would help determine
whether the method identifies biologically meaningful spatial structure beyond smoothing,
increased feature dimensionality, or clustering.
\end{rev}

\noindent This was the most important comment in the review, and we have addressed it with
three new controls, chosen from the reviewer's menu so that together they isolate the three
mechanisms the reviewer names. All three are reported in a new subsection of the Results
and a new supplementary figure, and the code and summary outputs are in the public
repository.

Throughout, spatial coherence of the resulting domain map is quantified with the
entropy-based local indicator of spatial association (ELSA) computed on the categorical domain
labels, at neighbourhood radii of 50 and 100 units. ELSA is negatively oriented: lower
values mean more spatially coherent domains. The domain-generating pipeline is held fixed
across observed and control runs, so any difference is attributable to the manipulation
alone.

\subsubsection*{Control A: ligand-receptor pairing shuffle}

We scrambled the assignment of receptors to ligands across mechanisms, holding the fitted
per-feature fields, the feature dimensionality, and the clustering procedure exactly fixed,
and repeated the convolution, clustering, and scoring for each of 99 scrambled pairings.

The observed pairing produced more spatially coherent domains than the scrambled pairings:
ELSA at radius 50 was 0.057 for the true pairing versus a null mean of 0.066, with the null
ranging from 0.053 to 0.078. Ninety-seven percent of scrambled pairings were less coherent
than the observed one, giving a one-sided permutation $p = 0.040$ at both radii.

Two aspects of this result deserve honest statement. First, the effect is modest, and it
should be: even a scrambled pairing recombines fields that are themselves genuinely
spatially structured, so this null is deliberately conservative and tests only the
incremental contribution of the biological pairing. Second, the number of domains does not
distinguish the observed pairing from the null at all (20 domains observed, against a null
range of 19 to 26, $p = 0.98$). The true pairing yields domains that are more spatially
coherent, not more numerous. We state both points in the manuscript.

\subsubsection*{Control B: coordinate shuffle with full model refitting}

This is the direct test of whether the spatial smoother manufactures coherent structure
from data with no spatial signal. We permuted the assignment of observed counts to spatial
locations, destroying spatial structure while preserving the count distribution exactly,
and then refit every feature's spatial model from scratch on the permuted data before
running the identical convolution, clustering, and scoring pipeline. This required
refitting all 912 features under each of 49 permutations, which we carried out on a
high-throughput computing cluster.

The result is unambiguous. The real data yielded 20 domains with ELSA at radius 50 of
0.075. The 49 coordinate-shuffled replicates yielded 4 to 12 domains with ELSA at radius 50
averaging 0.470, ranging from 0.412 to 0.583. Not one of the 49 shuffles came within a
factor of five of the observed coherence; the null mean is 6.3 times the observed value.
The permutation $p$-value is 0.020, which is the smallest value attainable with 49
permutations, so the test is limited by the number of replicates rather than by the
separation between the observed value and the null.

We regard this as the substantive answer to the reviewer's concern. Fitting a spatial
Gaussian process to spatially meaningless data does not produce coherent tissue domains; it
produces a small number of incoherent fragments. The domains we report therefore reflect
spatial structure present in the data, not structure imposed by the smoother.

\subsubsection*{Control C: concordance with independent anatomical annotation}

As a positive control we compared the emergent domains to the anatomical annotations
published with the source dataset by Chen et al.\ (2022) for the same tissue section. These
annotations were produced independently of our work, by a different method, so they
constitute external evidence rather than a circular check.

All 26{,}000 modeled spatial bins matched an annotated bin. Agreement was substantially
above chance: adjusted Rand index 0.401 and normalized mutual information 0.604, against a
permutation null whose mean adjusted Rand index was 0.00003 and whose maximum over 499
permutations was 0.001, giving $p = 0.002$. The mean size-weighted purity of the emergent
domains with respect to anatomical labels was 77.6\%. Individual domains recover recognized
structures cleanly: spinal cord at 97\% purity, liver at 96\%, one heart domain at 99\%,
meninges at 92\%, with lung, muscle, thymus, and cartilage primordium also recovered.

Again we state the result carefully. An adjusted Rand index of 0.401 is substantial but far
from perfect agreement, and the disagreement is structured rather than random: our
signaling domains subdivide several organs, most visibly the heart, into multiple
territories. We interpret this as the expected consequence of measuring a different
modality. Anatomical annotation partitions tissue by cell-type composition, whereas our
domains partition it by ligand-receptor signaling context, and the two need not coincide.
The appropriate conclusion is that the emergent domains recover major anatomical structures
while resolving additional substructure, not that they reproduce the anatomical
segmentation.

\subsubsection*{What the three controls establish together}

Control A shows that the biological ligand-receptor pairing contributes to the coherence of
the domains, modestly. Control B shows that the spatial structure is real rather than an
artifact of smoothing, decisively. Control C shows that the resulting structure corresponds
to recognized anatomy, while also resolving substructure within it. We believe these
address the reviewer's question of whether the method finds meaningful structure beyond
smoothing, dimensionality, and clustering.

We have not added simulated datasets with known ground truth, the remaining item on the
reviewer's menu. Our reasoning is that a simulation would have to be generated under a
spatial model, and any simulation we could write would be close enough to our own fitted
model that recovering the truth would demonstrate self-consistency rather than validity.
The coordinate-shuffle control tests the same concern using the real data and without that
circularity.

\subsection{R2.9. Reproducibility documentation} \status{PARTLY DONE}

\begin{rev}
Reproducibility also needs improvement. Although the repository contains a substantial
amount of code, the execution order of scripts, required packages and versions, required
input files, intermediate objects, figure-generating scripts, and complete workflow are not
sufficiently documented. An external researcher would likely need to reverse engineer the
analysis to reproduce the main results.
\end{rev}

\noindent This criticism was fair and we have acted on it. The repository now documents,
for each analysis, the execution order of the scripts, the required input files and where to
obtain the public ones, the intermediate objects each stage produces and consumes, and the
package versions used.

We have also drawn an explicit line between artifacts that are tracked in the repository and
large derived objects that are not. Every untracked derived object is now regenerated by a
tracked script, and the documentation names that script. Where a generating script did not
survive, we say so explicitly rather than implying the artifact is reproducible.

\todo{Only the new controls directory is documented to this standard so far. The main paper
pipeline under \texttt{Supplements for Publication/code/} still needs the same treatment:
run order, inputs, intermediates, figure scripts, and a session-info or lockfile capture.
This is the largest remaining task after the figures, and the reviewer will check it.}

\subsection{R2.10. Scientific positioning relative to existing work} \status{DRAFTED}

\begin{rev}
The scientific positioning of the manuscript should also be strengthened. The authors
should more clearly discuss related work in Gaussian Process-based spatial transcriptomics,
Bayesian spatial count modeling, spatial variance component analysis, and spatial
ligand-receptor inference. Relevant methods such as SpatialDE, BOOST-GP, Spatial Variance
Component Analysis, and SpatialDM should be discussed to clarify which aspects of the
manuscript are novel and which represent combinations, extensions, or reinterpretations of
existing approaches.
\end{rev}

\noindent We have added a positioning subsection that discusses each of these method
families and states plainly which parts of our contribution are new and which are
combinations or reinterpretations of established ideas.

In outline: SpatialDE introduced Gaussian-process testing for spatially variable genes, and
we share its use of a GP prior over space but not its goal, since we estimate and use the
field itself rather than testing features for spatial variability. BOOST-GP is the closest
methodological relative, in that it also places a Gaussian process within a Bayesian count
model for spatial transcriptomics; our contribution relative to it is not the count-GP
combination, which we do not claim as novel, but the downstream use of the posterior fields
as continuous objects that are convolved and then analyzed jointly. Spatial Variance
Component Analysis decomposes expression variance into spatial and non-spatial components,
which is a variance-partitioning aim rather than a field-estimation one. SpatialDM tests
ligand-receptor pairs for spatial association using a bivariate Moran's statistic, which
targets significance of individual interactions where we target a continuous inferred field
per interaction with quantified uncertainty.

We have also added a discussion of CytoSignal (Liu et al., \emph{Nature Genetics}, 2026),
concurrent work that appeared after our submission and that analyzes the same Stereo-seq
mouse embryo dataset. It infers spatially resolved ligand-receptor signaling at cellular
resolution using a score with a spatial-permutation null, and validates against experimental
proximity-ligation ground truth. We do not claim novelty for spatially resolved
ligand-receptor mapping in light of that work, and we have removed language that implied it.

Stating our contribution precisely, and more narrowly than the original submission did: we
contribute a generative Bayesian geostatistical treatment of spatialomics data with explicit
uncertainty quantification propagated through to the interaction fields, an argument for
modeling raw counts as absolute binding density rather than normalized expression, a
comparison of convolution operators showing that the common minimum outperforms the product
operator conventionally used, and the thesis that clustering the full high-dimensional
signaling field without dimensionality reduction produces tissue domains that correspond to
anatomy. The combination of a Gaussian process with a Bayesian count likelihood is not
itself new, and we now say so.

\todo{Write this into a new Discussion subsection and add the five citations. Verify the
characterization of each method against its paper rather than from summary knowledge,
particularly BOOST-GP and SVCA. Confirm the CytoSignal publication details and that the
concurrency framing is accurate with respect to both submission timelines.}

\subsection{R2.11. Generalizability claims} \status{DRAFTED}

\begin{rev}
Some claims of generalizability appear broader than the evidence provided. The study is
demonstrated using a selected region from a public Stereo-seq mouse embryogenesis dataset.
This is useful as a proof of concept, but spatial omics platforms differ in resolution,
capture efficiency, sparsity, noise structure, segmentation assumptions, field of view, and
biological scale. The authors should moderate broad claims about applicability to most
spatial omics datasets or provide additional analyses and practical guidance for different
platforms.
\end{rev}

\noindent We accept this and have moderated the claims rather than attempting to justify
them with additional platforms.

The revised manuscript states the scope of the demonstration explicitly wherever a general
claim previously appeared: one field of view, from one section, of one embryo, on one
platform, at one binning resolution. Claims of the form ``applicable to spatial omics data''
have been narrowed to describe what we demonstrate, with the extension to other platforms
stated as an expectation and its conditions given rather than asserted.

We have added a short paragraph on what would need to hold for the approach to transfer,
which we think is more useful to a reader than a broad claim: the observations must be
counts with an interpretable exposure or offset, the spatial support must be dense enough
that a Gaussian process is identifiable at the length scales of interest, and at
single-cell-resolution platforms the binning step we rely on would need to be replaced by an
explicit treatment of cell geometry. We also note that platforms with substantially sparser
capture may favour a zero-inflated likelihood, which our Table 1 comparison found
unnecessary at this sparsity but which the framework accommodates.

\todo{Apply the softening pass across the manuscript. The Abstract, the end of the
Introduction, and the Discussion are where the broad claims live. Check that the moderated
scope is consistent everywhere, including the Abstract.}

\subsection{R2.12. Caution in biological interpretation} \status{DRAFTED}

\begin{rev}
Finally, the biological interpretation should remain cautious. I appreciate that the
authors acknowledge that transcript abundance is only used as a proxy for protein abundance.
However, because downstream analyses involve inferred ligand-receptor interactions and
morphogenic signaling fields, the manuscript should avoid language that implies direct
evidence of protein-level interactions, ligand-receptor binding, signaling activity, or
information transmission. Terms such as ``proxy,'' ``inferred field,'' ``model-derived
estimate,'' or ``potential interaction'' would be more appropriate.
\end{rev}

\noindent We agree, and we have revised the language throughout.

We have replaced formulations that implied observed protein-level events with formulations
describing model-derived estimates from transcript measurements. Descriptions of
ligand-receptor quantities now consistently use ``inferred interaction field'' or
``potential interaction'' rather than wording that implies binding was observed; references
to signaling activity are now phrased as transcriptional evidence consistent with potential
signaling; and language implying information transmission between regions has been removed
or explicitly marked as interpretation.

We have also added an explicit statement of the inferential chain and its assumptions in the
Discussion, so a reader can see exactly how far the claims are from the measurement:
transcript counts are a proxy for transcript abundance, transcript abundance is a proxy for
protein abundance, co-localization of ligand and receptor proxies is a proxy for the
opportunity for interaction, and none of these steps is directly observed here.

\todo{Perform the terminology pass across the full manuscript. Read the Results and
Discussion specifically for verbs implying observation of binding or signaling. Keep
terminology consistent: pick one term per concept and use it throughout.}

\section{Additional corrections}

\ifdraftmode
\noindent\textcolor{red}{\textbf{Internal note:}} the item below was found by us, not by a
reviewer. It should still be disclosed in the letter, since the reviewer will notice the
numbers changed.
\medskip
\fi

\noindent In preparing this revision we identified an internal inconsistency in the
original manuscript's reporting of the number of features and interactions analyzed, which
we have corrected throughout. The curated feature list contains 918 features; of these, 912
produced converged model fits and were carried forward. The FANTOM5 mechanism list yields
1477 ligand-receptor interactions for which both partners appear among the 912 converged
features, not the 1481 reported in the original text and in the caption to Table 1. The
caption to Table 1 also described 1481 ``modeled transcriptomic features,'' conflating
interactions with features. All counts are now stated consistently, and we have added a
sentence explaining the attrition from 918 to 912.

\todo{Verify these four numbers against the canonical run before this goes in, and correct
every occurrence in the manuscript and figure captions rather than only the ones we
noticed.}

\end{document}
